% Peekaboo-VL report: 08_discussion. Included from ../peekaboo_report.tex.
\section{Discussion and Limitations}

\textbf{What the video adds.} On ReVA, a text-only model already
reaches 56\% and a 9B video model without fine-tuning 80\%. So most of the remaining
room is in a few tasks: Temporal Grounding, Change Detection and Geometric
Relation. These are exactly the tasks the blind scorer cannot solve. We
suggest that results on ReVA be reported together with a trained blind
baseline and per-task numbers, so readers can see what comes from the
video.

\textbf{Telling instead of learning.} Text hints are cheap, easy to read
and check, and work with any model. But they have two risks that learned
motion tokens do not. First, their wording interacts with the answer
options, as Section~\ref{sec:cue} shows. Second, a wrong hint is stated
just as confidently as a right one. Fine-tuning lets the model learn when
to ignore the hint, and our test with wrong hints shows it does: it follows
the hint on camera-motion questions and ignores it elsewhere.

\textbf{Takeaways we are carrying forward.}
\begin{enumerate}
\item Check how the test relates to the training data before building
anything. Here, the test asks new questions about seen videos, and the
validation set is mostly copies.
\item Build a dev set that copies the test setting, and weight it by the
test's task mix.
\item Train a blind baseline first. It tells you what the video adds and
which tasks need it.
\item When you add text hints, use the words of the answers.
\item Shuffle answer options during training; most letter bias then
disappears, but cheap debiasing (PriDe) still won on the test.
\item Do not trust dev gains of about one point, whether from decoding,
hints or model combinations; several of ours vanished on the test.
\item Do not change the frame layout at test time unless the model was
trained with that layout.
\end{enumerate}

\textbf{Limitations.} Most ablations ran on our first model (57\% of the
fit questions, 16 frames) or on the packed recipe, not on the final model,
so their deltas are directions for the final recipe, not measurements of
it. Our dev split shares clips with the fit set on purpose: it measures
the ``seen video, new question'' setting of this challenge, not how well
the model handles new videos. A fifth of dev are test questions whose
answers the official validation file contains, which makes dev a closer
proxy for the test than a disjoint split would be; we never used those
answers to pick test predictions. We did not use the
training answers about a test clip as context at test time (``video
memory''); the organizers confirmed that this is not allowed, although
training on those labels is. Finally, the probe's four-number motion model
cannot represent camera turns on strongly tilted views; a homography, as
used in our object-motion hint, is the natural next step. Our headline
comparison is against ReMoSense's published numbers rather than a
reproduction, and every configuration is a single seed; both are about our
budget, and both are reasons to read the differences as directions rather
than exact gaps.

\textbf{Broader impact and ethics.} This work studies aerial video
understanding. Drone video is useful for disaster
response, agriculture and traffic monitoring, and it also enables surveillance of
people who did not consent to being recorded. The ReVA dataset is released
for non-commercial research only and may not be redistributed, and we follow
those terms and release no new video data. Our most reusable artifact is a
few-millisecond CPU motion probe that runs on any video. The broader
caution is methodological: when a benchmark's text alone is over half
solvable and its test videos appear in training, reported accuracy can
overstate genuine visual understanding. A recent survey of aerial video
benchmarks~\cite{angelidis2026aerialgap} reached the same conclusion from
the other direction: none of the question-answering benchmarks it audited
reports a blind floor.

